\subsection{Threshold Sensitivity Within the Selected Configuration}
Table~\ref{tab:threshold} holds the RF\_leaf\_4 candidate fixed and varies only the decision threshold using its saved out-of-fold probabilities. This comparison isolates a decision-rule effect that the V1--V2 comparison cannot isolate. Moving from 0.50 to 0.40 raises mean accuracy from 85.96\% to 90.37\% and BA from 81.58\% to 83.66\%. Resistant recall increases by 5.45 percentage points, whereas susceptible recall decreases by 1.28 points.

\begin{table}[t]
\caption{Threshold Sensitivity for the Selected Forest (Percent)}
\label{tab:threshold}
\centering\footnotesize
\begin{tabular}{crrrrr}\toprule
$\tau$ & Acc. & BA & S recall & R recall & Worst BA\\\midrule
0.30&90.17&80.47&65.59&95.35&60.29\\
0.40&90.37&83.66&73.35&93.98&67.15\\
0.50&85.96&81.58&74.63&88.52&68.24\\
0.60&80.30&78.25&75.51&80.99&65.32\\\bottomrule
\end{tabular}
\vspace{3pt}
\parbox{\columnwidth}{\footnotesize All values use the same selected-candidate development predictions. Worst BA is the lowest fold balanced accuracy. These are threshold-selection results, not an external threshold validation.}
\end{table}

Reducing the threshold further to 0.30 gives slightly lower accuracy than 0.40 but a much larger reduction in susceptible recall. Raising it to 0.60 recovers little susceptible recall while losing resistant recall. This pattern explains why the recorded development objective selects 0.40 from the grid. It does not establish that 0.40 is optimal for a new hospital, laboratory, or prevalence setting.

The best mean BA and best worst-fold BA are also not identical objectives. Among the displayed thresholds, 0.50 has a higher worst-fold BA than 0.40 even though its mean BA is lower. The implemented tie breaker cannot override a difference in the primary mean objective. If a future application prioritizes minimum subgroup performance, that requirement should be stated before model selection and evaluated independently, rather than introduced after inspecting a difficult fold.

\subsection{Pooled Confusion Counts and Project Concentration}
Pooling the selected model's out-of-fold predictions gives 63 true susceptible predictions, 20 susceptible isolates predicted resistant, 22 resistant isolates predicted susceptible, and 340 true resistant predictions. Table~\ref{tab:confusion} presents the counts without averaging. The resulting pooled accuracy is $403/445=90.56\%$, pooled susceptible recall is $63/83=75.90\%$, pooled resistant recall is $340/362=93.92\%$, and pooled BA is 84.91\%. These pooled values coexist with, and do not replace, the reported mean-fold scores.

\begin{table}[t]
\caption{Pooled V2 Out-of-Fold Confusion Matrix}
\label{tab:confusion}
\centering\small
\begin{tabular}{lrrr}\toprule
Actual phenotype & Predicted S & Predicted R & Total\\\midrule
Susceptible &63&20&83\\
Resistant &22&340&362\\\midrule
Total &85&360&445\\\bottomrule
\end{tabular}
\end{table}

The difference between pooled and mean-fold BA is driven by weighting. Fold 3 contains 26 susceptible isolates and has the poorest susceptible recall; the fold mean gives that fold the same overall weight as the others. Pooled class recall instead weights each susceptible isolate equally. Neither calculation removes the failure. Reporting the formula, denominator, and aggregation rule prevents a higher aggregate from being selected merely because it is more favorable.

Joining the V2 predictions to the saved BioProject manifest shows that errors are concentrated in a few projects. PRJNA278886 contributes 20 of the 42 total errors: 16 false-resistant and four false-susceptible predictions among 109 isolates. PRJNA376414 contributes 14 errors, including 12 false-susceptible predictions among 101 isolates. Together these projects contain 210 of 445 isolates and account for 34 of 42 errors. Their error profiles differ, so one overall explanation for the model's mistakes would be premature.

\begin{table}[t]
\caption{Largest Error-Contributing Projects in V2}
\label{tab:projects}
\centering\footnotesize
\begin{tabular}{lrrrrr}\toprule
BioProject & $n$ & S & R & S as R & R as S\\\midrule
PRJNA278886&109&22&87&16&4\\
PRJNA376414&101&10&91&2&12\\
PRJNA288601&87&2&85&1&2\\
PRJNA396774&7&5&2&1&1\\\bottomrule
\end{tabular}
\vspace{3pt}
\parbox{\columnwidth}{\footnotesize Projects are shown in descending total error count. Counts come from selected-model development predictions, not the historical V1 test.}
\end{table}

Table~\ref{tab:projects} also illustrates why project-level percentages can be unstable. PRJNA288601 has only two susceptible isolates, so one susceptible error changes that class recall by 50 percentage points. PRJNA396774 contains seven isolates in total. Such small denominators warrant displaying counts alongside any percentage. These descriptive project comparisons identify where evidence review is needed; they do not estimate an institution's population resistance prevalence or establish that project membership causes an error.

\subsection{Interpretation of Regenerated Sequence Evidence}
The processing checkpoint supports a limited but useful conclusion: the workflow can connect public accessions to nucleotide assemblies and reproducible annotation outputs for a substantial part of the development cohort. It cannot yet answer whether regenerated evidence improves classification. The completed assembly subset and remaining read-derived subset have different data-availability conditions. Training only on the convenient completed subset could change the cohort and make a feature comparison inseparable from a sampling comparison.

Completeness should therefore be evaluated at the manifest level. An isolate with a failed download is different from an isolate with an assembly warning, and both differ from an isolate with a valid annotation containing no retained resistance feature. Treating all three as an all-zero vector would hide distinct evidence states. The saved processing records support keeping these states separate while unresolved cases are reviewed. No unresolved sample should be removed because its current model prediction is wrong.

The mapping checks also require contextual interpretation. Aligning reads back to their own assembly tests whether assembled sequence is supported by those reads, but missing sequence can remain absent from both the reference and the evaluated breadth denominator. Similarly, a close species match is compatible with additional contamination that a dedicated screen might detect. The retained uniform-quality flag demonstrates why quality summaries should not be reduced to a single pass/fail score. These are reasons to retain raw reports, not evidence that a specific isolate is biologically invalid.

\subsection{Requirements for the Next Evaluation}
The next evaluation should distinguish a feature-development stage from a final validation stage. First, the complete sequence-derived representation should be defined, including partial calls, mutation states, missingness, database version, and quality-review rules. Candidate comparison can then proceed using development data with the overlap restrictions applied to that representation. Only after the pipeline and threshold are frozen should an independent cohort be used to estimate generalization.

The independent cohort should document the phenotype testing method and interpretive standard, retain accession provenance, and include enough susceptible and resistant isolates to interpret both error denominators. Evaluation should report accuracy, BA, class recalls, ranking performance, and confusion counts under the frozen decision rule. Any calibration assessment or alternative operating threshold should be declared as a separate analysis. These are requirements for a future study; they are not completed experiments in the present paper.

Operational surveillance would require an additional evaluation beyond isolate classification. For example, measuring whether flagged isolates support earlier investigation requires timestamps, a reference definition of an event, and an explicit alert policy. None of those outcomes is available in the present saved experiment. The immediate result is therefore a reproducible prediction-development record and a sequence-processing feasibility demonstration. Keeping that boundary explicit makes later claims testable when operational data become available.

